\documentclass[10pt,a4paper]{amsart}
\usepackage[margin=19mm]{geometry}
\usepackage{amsmath,amssymb,mathtools}
\usepackage[scr=rsfs,cal=euler]{mathalfa}
\usepackage{xfrac}
\usepackage[unicode,hidelinks,bookmarks=false]{hyperref}
\newcommand{\ie}{i.e.\ }
\newcommand{\cf}{cf.\ }
\newcommand{\resp}{respectively\ }
\newcommand{\Subsection}{\subsection}
\providecommand{\nobreakdash}{\nobreak\hbox{-}}
\setlength{\emergencystretch}{2em}
\multlinegap0pt
\newtheorem{proposition}{Proposition}
\newcommand{\R}{\mathbb{R}}
\newcommand{\e}{\mathrm{e}}
\newcommand{\Id}{\mathrm{Id}}
\newcommand{\supp}{\mathrm{supp}}
\newcommand{\g}{\mathbf{g}}
\renewcommand{\i}{\mathrm{i}}
\renewcommand{\d}{\mathrm{d}}
\title[Sharp lower bound for $\sin(t\sqrt{L})/\sqrt{L}$ on $ax+b$ groups (Proposition 5.1)]{Sharp $L^p$-estimates for wave equation on $ax+b$ groups: the sharp polynomial-in-time lower bound (Proposition 5.1)}
\author{Yunxiang Wang, Lixin Yan}
\date{Source: arXiv:2506.17531v2 (CC BY 4.0)}
\begin{document}
\begin{abstract}
Let $G$ be the group $\mathbb{R}_+\ltimes \mathbb{R}^n$ endowed with Riemannian symmetric space metric $d$ and the right Haar measure $\mathrm{d} \rho$ which is of $ax+b$ type, and $L$ be the positive definite left-invariant distinguished Laplacian on $G$. Let $u=u(t,\cdot)$ be the solution to $\partial_t^2u+Lu=0$ with initial conditions $u|_{t=0}=f_0$ and $\partial_tu|_{t=0}=f_1$. In this article we show that for a fixed $t\in\mathbb{R}\setminus\{0\}$ and every $1<p<\infty$, 
\begin{align*}
\|u(t,\cdot)\|_{L^p(\mathrm{d}\rho)}\leq C_p\,\Big( (1+|t|)^{2|1/p-1/2|}\|f_0\|_{L^p_{\alpha_0}(\mathrm{d}\rho)}+(1+|t|)\,\|f_1\|_{L^p_{\alpha_1}(\mathrm{d}\rho)}\Big) 
\end{align*}
if and only if
\begin{align*}
\alpha_0\geq n\left|{1\over p}-{1\over2}\right| \quad \mbox{and}
\quad \alpha_1\geq n\left|{1\over p}-{1\over2}\right|-1.
\end{align*} 
The polynomial-in-time growth $1+|t|$ multiplying $\|f_1\|_{L^p_{\alpha_1}(\mathrm{d}\rho)}$ in the above estimate is best possible. 
This endpoint result settles the conjecture raised by M\"{u}ller and Thiele [Studia Math. \textbf{179} (2007)].


\end{abstract}
\maketitle

\noindent\emph{Source and licence.} Sharp $L^p$-estimates for wave equation on $ax+b$ groups. Version arXiv:2506.17531v2 (CC BY 4.0), distributed under the Creative Commons Attribution 4.0 International licence, as recorded in the arXiv licence metadata for this exact version and shown on the abstract page https://arxiv.org/abs/2506.17531v2. Full rights notice: licenses/arxiv-2506.17531-CC-BY-4.0.txt.\par\smallskip

\noindent\textit{Editorial scope.} Counted unit: the proposition printed as Proposition 5.1 of the article (source0.tex lines 989--994, source label \texttt{prop5.1}), delivered together with the article's own complete original proof of it (source0.tex lines 995--1088, one \texttt{proof} environment of 94 lines that proves the lower bound with no gap and no deferral). No mathematics is written, repaired or re-derived: statement and proof are the article's own text line for line, in the article's own notation ($G=\R_+\ltimes\R^n$, $L$, $p$, $\alpha$, $t$, $q_t$, $S_t$, $\mathcal T_t$, $\mathcal N_t(S_t)$, $\rho$). Required context retained uncounted: the article's abstract (lines 73--83), which states the sharpness assertion that the counted proposition strengthens; the Cauchy problem (e1.4) for the wave equation, whose unit-speed propagation the proof uses (lines 122--125); and the left-invariant Riemannian metric (riem) with which the article defines the geometry of $G$ (lines 205--208). Every definition, numbered display and label of those lines is retained unchanged. Omitted from this package and not redistributed: the article's preamble and title page (lines 1--70); the end of its front matter (lines 84--96); its whole introduction apart from the retained Cauchy problem, including the statements of its main theorem and of the Mueller--Thiele conjecture it resolves (lines 97--121, 126--203); its second section on the hyperbolic space, harmonic analysis and Hardy spaces apart from the retained metric (lines 209--430); its third section on the explicit asymptotic expansions of spherical functions apart from nothing retained inside it (lines 431--608); its fourth section, the whole sufficiency proof of its main theorem with the Littlewood--Paley and multiplier propositions (lines 609--880); its fifth section, the whole necessity proof of its main theorem with the sharpness propositions that lead up to the counted one (lines 881--988); everything after the counted proof, namely the remaining sharpness corollaries and the closing remarks (lines 1089--1104); and the article's bibliography (lines 1105--1156) together with its closing end-document line (line 1157). Nothing inside a retained excerpt is omitted or altered: every retained body is delivered exactly as the article's lines are written, so no omission receipt of any kind accompanies this package. Disclosed dependence on omitted material: the counted proof uses results of the article's omitted earlier sections through three references to the published literature that it names at the point of use - the multiplier estimates of Mueller and Thiele, the unit-speed propagation theorem of Sogge and a classical Fourier-transform estimate of Stein - and it uses the symbol class $S^\sigma$ and the spectral calculus for $L$ that the omitted sections set up. Citations: the retained excerpts carry exactly four citation commands, all inside the counted proof, and together they cite exactly three keys of the article's own bibliography, \texttt{MuTh}, \texttt{Sog14} and \texttt{St}. Those three entries are transcribed verbatim from the article's own in-source bibliography (lines 1105--1156, 25 entries) with the article's own printed labels, so the delivered citation marks read [14], [20] and [22] as the article's do. No other entry of the article's bibliography is transcribed or redistributed. Reference adaptation: none was needed and none was made. Every reference inside the delivered unit resolves inside it: the label e1.4 of the retained Cauchy problem, the label riem of the retained Riemannian metric, and the labels riem-sy, eq:6.2, eq:prop5.1claim, eq:prop5.1final and e6.3.0 of the numbered displays inside the counted proof, together with the label prop5.1 of the counted statement itself. No unresolved reference or citation is produced. Printed slips kept literal: no printed word, symbol, hypothesis or number of the counted statement or of its proof was altered; in particular the retained text keeps the article's own naming of the sets $\mathcal T_t$ and $S_t$ and its own use of the ramified notation $\e^{-2s}\d y^2$. Layout and class: the article's own class is \texttt{amsart} with the option 12pt and a4paper, and it loads amssymb, amsxtra, paralist, mathrsfs, upgreek, xcolor, ulem and hyperref; no publisher class or style file is used or shipped, so none travels with this package and none is redistributed. The wrapper is the lane's pinned \texttt{amsart} editorial wrapper at 10pt on A4, which adds the article's own macro definitions that the retained text needs ($\R$, $\e$, $\Id$, $\supp$, $\g$, $\i$, $\d$) and declares the single theorem-like environment the retained text needs (\texttt{proposition}). The delivered numbering is the unit's own: the counted Proposition 5.1 prints as Proposition 1. The article's 12pt measure differs from the wrapper's 10pt, and the article's magenta/cyan hyperlink colours are not reproduced by the wrapper, which uses black links; those are page-appearance differences of the wrapper only, and no formula, symbol, hypothesis or argument changes. Assets: no retained span contains a figure environment, a graphics-inclusion command, a PDF-page inclusion, a table or a TikZ picture, no image file is copied into this package, the package declares no asset dependency, and no image file is redistributed with it. Licence: version arXiv:2506.17531v2 of this article is distributed under the Creative Commons Attribution 4.0 International licence, as recorded in the arXiv licence metadata for this exact version and shown on the abstract page https://arxiv.org/abs/2506.17531v2. A full rights notice is delivered at licenses/arxiv-2506.17531-CC-BY-4.0.txt. Characters: every retained excerpt is pure ASCII, so the delivered engine, XeLaTeX with its default Latin Modern text font, reports no missing character. No glyph is substituted and no word is rewritten.
Let $u=u(t,(x,y))$ be the solution to the Cauchy problem
\begin{align}\label{e1.4}
{\partial^2u\over\partial t^2}+Lu=0,\quad u|_{t=0}=f_0,\quad \left.{\partial u\over\partial t}\right|_{t=0}=f_1.
\end{align}

The space $G$ endowed with the left-invariant Riemannian metric given by
\begin{align}\label{riem}
\g=x^{-2}(\d x^2+\d y^2)
\end{align}

\begin{proposition}\label{prop5.1}
Suppose $1\leq p\leq\infty$ and $t>0$. Then for all $\alpha\in\R$
\begin{align*}
\left\|{\sin(t\sqrt{L})\over\sqrt{L}}\right\|_{L^p_\alpha(G)\to L^p(G)}\geq Ct.
\end{align*}
\end{proposition}

\begin{proof}
To simplify calculations, we let $s=\log x$ and adopt $(s,y)$-coordinates. This time
\begin{align*}
\d\rho(s,y)=\d s\d y,\qquad L=-X^2-\sum_{j=1}^nY_j^2=-\partial_s^2-\e^{2s}\Delta_y
\end{align*}
for $\Delta_y$ being the standard Laplacian on $\R^n$, while the Riemannian metric \eqref{riem} becomes
\begin{align}\label{riem-sy}
\g=\d s^2+\e^{-2s}\d y^2.
\end{align}

We let $\eta\in C^\infty_c([-2,2])$ with $0\leq \eta\leq 1$ and $\eta\equiv 1$ on $[-1,1]$, and define
\begin{align*}
q_t(s,y):=\eta\!\left({s\over 2t+1}+3\right)\prod_{j=1}^n\eta(y_j).
\end{align*}
Then
\begin{align*}
\supp q_t\subset \widetilde{\mathcal{T}}_t:=\{-10t-5\leq s\leq -2t-1,\, |y_j|\leq 2\}
\end{align*}
and
\begin{align*}
q_t\equiv 1\quad\text{on }\mathcal{T}_t:=\{-8t-4\leq s\leq -4t-2,\, |y_j|\leq 1\}.
\end{align*}
\smallskip

We first claim that for $\widetilde{\alpha}\in\R$,
\begin{align}\label{eq:prop5.1claim}
\|(\Id+L)^{\widetilde{\alpha}}(q_t)\|_{L^p(G)}\leq C\, (1+t)^{1/p}.
\end{align}
To see this, let $m=\max\{\lceil\widetilde{\alpha}\rceil,0\}$. Then
\begin{align}\label{eq:6.2}
\|(\Id+L)^{\widetilde{\alpha}}(q_t)\|_{L^p(G)}\leq \underbrace{\|(\Id+L)^{\widetilde{\alpha}-m}\|_{L^p(G)\to L^p(G)}}_{\mathbf I}\underbrace{\|(\Id+L)^m(q_t)\|_{L^p(G)}}_{\mathbf J}.
\end{align}
Since $(\Id+L)^m$ is a linear combination of the following differential operators:
\begin{align*}
\e^{2ks}\partial_s^\ell\Delta_y^k,\qquad 2k+\ell\leq 2m,
\end{align*}
and on $\supp q_t\subset \widetilde{\mathcal{T}}_t$
\begin{align*}
\e^{2ks}\partial_s^\ell\Delta_y^k(q_t)\leq C_{k,\ell},
\end{align*}
we have
\begin{align*}
\mathbf J\leq C_m\,(\rho(\widetilde{\mathcal{T}}_t))^{1/p}\leq C\, (1+t)^{1/p}.
\end{align*}

On the other hand, if $\widetilde{\alpha}=0,1,2,\dots$ we have $\mathbf I=1$; otherwise by \cite[Propositions 4.1, 5.2 and 5.7]{MuTh} and classical estimates on Fourier transforms of symbols \cite[p. 241]{St}, $(\Id+L)^{\widetilde{\alpha}-m}$ has an $L^1$-integrable convolution kernel, whence $\mathbf I\leq C$. In view of \eqref{eq:6.2}, we obtain \eqref{eq:prop5.1claim}.
\smallskip

With \eqref{eq:prop5.1claim}, to show Proposition \ref{prop5.1} it suffices to establish
\begin{align}\label{eq:prop5.1final}
\left\|{\sin(t\sqrt{L})\over\sqrt{L}}(q_t)\right\|_{L^p(G)}\geq Ct\,(1+t)^{1/p}.
\end{align}
To see this, by the argument showing \cite[Theorem 2.4.2]{Sog14}, we know that the wave equation \eqref{e1.4} has \textit{unit speed propagation}. In other words, for
$$
(s,y)\in S_t:=\{-7t-4\leq s\leq -5t-2,\, |y_j|\leq 1/2\},
$$
the functions
$$
{\sin(t\sqrt{L})\over\sqrt{L}}(q_t)(s,y)
$$
depend only on $q_t$ restricted to the $t$-neighborhood $\mathcal N_t(S_t)$ of $S_t$.

We claim that $\mathcal N_t(S_t) \subset \mathcal T_t$. Indeed, suppose $(s,y)\in S_t$, and let $(s(\tau),y(\tau))$ be any arc-lengthed geodesic ray starting from $(s,y)$. By \eqref{riem-sy}, this means that
\begin{align*}
|s'(\tau)|^2+\e^{-2s(\tau)}\sum_{j=1}^n|y_j'(\tau)|^2\equiv 1.
\end{align*}
Then for $0\leq t_0\leq t$,
\begin{align}\label{e6.3.0}
|s(t_0)-s|\leq \int_0^{t_0}|s'(\tau)|\,\d\tau\leq t,
\end{align}
and in view of \eqref{e6.3.0} and the fact that $s\leq -5t-2$,
\begin{align*}
|y_j(t_0)-y_j|\leq \int_0^{t_0} |y_j'(\tau)|\,\d \tau\leq t\sup_{0\leq\tau\leq t}\e^{s(\tau)}\leq t\,\e^{-4t-2}\leq {\e^{-3}\over 4}.
\end{align*}
Hence for $0\leq \tau\leq t$,
\begin{align*}
-8t-4\leq s(\tau)\leq -4t-2,\qquad |y_j(\tau)|\leq 1,
\end{align*}
whence $(s(\tau),y(\tau))\in \mathcal{T}_t$. The claim is verified.

To show \eqref{eq:prop5.1final}, we note that $u(t,\cdot)\equiv t$ is the solution to the Cauchy problem
\begin{align*}
{\partial^2u\over\partial t^2}+Lu=0,\quad u|_{t=0}\equiv 0,\quad\left.{\partial u\over\partial t}\right|_{t=0}\equiv1.
\end{align*}
As a consequence, by the unit speed propagation property and uniqueness of the solution (cf. \cite[pp. 12--13]{Sog14}), one has
\begin{align*}
{\sin(t\sqrt{L})\over \sqrt{L}}(q_t)\equiv t \qquad\text{on }S_t.
\end{align*}
Hence
\begin{align*}
\left\|{\sin(t\sqrt{L})\over \sqrt{L}}(q_t)\right\|_{L^p(G)}\geq {t\over 2}\,(\rho(S_t))^{1/p}\geq Ct\,(1+t)^{1/p}.
\end{align*}
This proves \eqref{eq:prop5.1final} and hence Proposition \ref{prop5.1}.
\end{proof}

\begin{thebibliography}{99}
\bibitem[14]{MuTh} D. M\"{u}ller and C. Thiele, Wave equation and multiplier estimates on $ax+b$ groups. \textit{Studia Math.} \textbf{179} (2007), no. 2, 117--148.
\bibitem[20]{Sog14} C.D. Sogge, \textit{Hangzhou lectures on eigenfunctions of the Laplacian}. Ann. of Math. Stud. \textbf{188}, Princeton University Press, Princeton, NJ, 2014. xii+193 pp.
\bibitem[22]{St} E.M. Stein, \textit{Harmonic analysis: Real variable methods, orthogonality and oscillatory integrals}. Princeton Math. Ser. \textbf{43}. Monographs in Harmonic Analysis \textbf{III}. Princeton Univ. Press, Princeton, NJ, 1993.
\end{thebibliography}
\end{document}
